\documentclass[12pt]{article}

% Packages
\usepackage[
  paperwidth=8.5in,
  paperheight=11in,
  margin=0.5in
]{geometry}
\usepackage{amsmath, amssymb}
\usepackage{booktabs}
\usepackage{tabularx}
\usepackage{graphicx}
\usepackage{hyperref}


\newcommand{\tacred}{FS-TACRED}
\newcommand{\qwensmall}{Qwen3-4B}
\newcommand{\gemmasmall}{Gemma3-4B}

\begin{document}

\paragraph{Algorithm of the new approach on top of the Error Taxonomy}
\begin{enumerate}
    \item Run initial prompt on validation set and collect errors.
    \item Build error taxonomy and group errors into categories.
    \item Select important categories based on frequency and coverage.
    \item Split categories into multiple overlapping clusters.
    \item For each cluster (having a category set):
    \begin{itemize}
        \item Generate few prompts by updating the initial prompt and keep the best one.
    \end{itemize}
    \item Select top prompts (each prompt comes from different clusters) based on their evaluation score and use them for inference.
\end{enumerate}

\paragraph{Experiment Settings of Baseline (Error Taxonomy):}
\begin{itemize}
    \item Select top 7 categories or enough to reach 70\% coverage of error counts.
    \item Generate 15 prompts.
    \item Select best 3 prompts for inference.
\end{itemize}

\paragraph{Experiment Settings of Clustered Approach:}
\begin{itemize}
    \item Use same selected categories by the baseline above.
    \item Divide the categories into 5 overlapping clusters.
    \item Generate 3 prompts per cluster (total = 15).
    \item Keep best 1 prompt per cluster.
    \item Select best 3 cluster prompts for inference.
\end{itemize}

\paragraph{Additional Details:}
\begin{itemize}
    \item Both methods generate 15 prompts (fair comparison).
    \item Each experiment is run 3 times and averaged the metrics.
    \item Also tested with broader setting: top 10 categories or up to 90\% coverage.
    \item Model usage: Qwen3-14B for error category generation and Qwen3-4B for inference.
\end{itemize}

\paragraph{Results} The clustered approach shows a small improvement based on the avg. of 3 runs.
For the default setting, F1 improves from 31.9 to 33.28.
For the broader setting (more errors used), F1 improves from 28.1 to 28.21.

% \bibliographystyle{plain}  
% \bibliography{custom}    

\end{document}
